% Blueprint: Summability of the zeros of the Riemann ξ-function
%
%   ∑_ρ |ρ|^{-(1+ε)} < ∞      (ε > 0),     in particular   ∑_ρ |ρ|^{-2} < ∞,
%
% the sum running over the zeros ρ of ξ counted with multiplicity.
%
% References: Titchmarsh, *The theory of the Riemann zeta-function* (2nd ed.) §3.6–§3.9
% \cite{TitchmarshZeta}; Davenport, *Multiplicative Number Theory* (3rd ed.) §11
% \cite{Davenport}.

\chapter{Summability of the zeros of the Riemann \texorpdfstring{$\xi$}{xi}-function}
\label{cha:zero_summability}

Throughout, $\xi$ is \texttt{riemannXi}, entire
(\texttt{differentiable\_riemannXi}) with $\xi(0)=\tfrac12\neq0$
(\texttt{riemannXi\_zero}). We write
\[
  D \;:=\; \texttt{MeromorphicOn.divisor riemannXi Set.univ}
  \;:\; \texttt{locallyFinsuppWithin Set.univ}\ \mathbb Z,
\]
the divisor of $\xi$. Because $\xi$ is analytic on all of $\mathbb C$, the divisor is
non-negative, $D\ge0$ (\texttt{AnalyticOnNhd.divisor\_nonneg}); $D(z)$ is the
multiplicity (order of vanishing, \texttt{meromorphicOrderAt riemannXi z}, an integer
since $\xi$ is holomorphic) of $\xi$ at $z$, $D(0)=0$ (\texttt{riemannXi\_zero} gives
$\xi(0)\neq0$), and on every compact set $D$ has finite support
(\texttt{locallyFinsuppWithin.finiteSupport}, the analytic zeros being isolated). For a
real radius $R$ we abbreviate by $D_R$ the restriction
$\texttt{MeromorphicOn.divisor riemannXi (Metric.closedBall 0\ R)}$, which equals
$\texttt{locallyFinsuppWithin.toClosedBall}\ R\ D$ (\texttt{toClosedBall\_divisor}); its
support lies in $\{z:|z|\le R\}$.

The single analytic input is the order-one growth bound proved in \Cref{cha:xi_order}:
there is a constant $C_\xi>0$ with $\|\xi(s)\|\le\exp(C_\xi\,|s|\log|s|)$ for $|s|\ge2$
(\Cref{thm:xo_order}, \texttt{norm\_riemannXi\_le\_exp}). Equivalently, on
the circle $|s|=R$ ($R\ge2$) one has the pointwise upper bound $\log\|\xi(s)\|\le
C_\xi\,R\log R$, finite where $\xi(s)\neq0$, and trivially $\le C_\xi R\log R$ at a zero
of $\xi$ on the circle, where mathlib's convention $\log 0=0\le C_\xi R\log R$ applies.

\section{The integrated counting function}
\label{sec:zs_jensen}

Mathlib's logarithmic counting function of a divisor is
\[
  N(R)\;:=\;\texttt{Function.locallyFinsuppWithin.logCounting}\ D\ R
  \;=\;\sum_{z}^{\flat} D_R(z)\,\log\!\frac{R}{|z|}\;+\;D(0)\,\log R,
\]
where $\sum^{\flat}$ is the \texttt{finsum} over the (finite, for fixed $R$) support of
$D_R$. As $D(0)=0$, the second term vanishes.

\begin{lemma}[Jensen bound on the integrated counting function]
\label{lem:zs_jensen}
\lean{logCounting_riemannXi_le}
\leanok
\uses{thm:xo_order}
There is a constant $C_N>0$ such that for all $R\ge2$,
\[
  N(R)\;=\;\texttt{logCounting}\ D\ R\;\le\;C_N\,R\log R .
\]
\end{lemma}

\begin{proof}
\leanok
\uses{thm:xo_order}
By Jensen's formula in the counting-function form
(\texttt{logCounting\_divisor\_eq\_circleAverage\_sub\_const}, valid since $\xi$ is
meromorphic on $\mathbb C$ and $R\neq0$),
\[
  N(R)\;=\;\operatorname{circleAverage}(\log\|\xi(\cdot)\|)\,0\,R
          \;-\;\log\|\texttt{meromorphicTrailingCoeffAt}\ \xi\ 0\|.
\]
Since $\xi$ is analytic at $0$ with $\xi(0)=\tfrac12\neq0$, its trailing coefficient at
$0$ is $\xi(0)=\tfrac12$ (\texttt{AnalyticAt.meromorphicTrailingCoeffAt\_of\_ne\_zero}),
so the constant term is $-\log\tfrac12=\log2$. For the circle average: the integrand
$\log\|\xi(\cdot)\|$ is circle-integrable (\texttt{circleIntegrable\_log\_norm\_meromorphicOn},
$\xi$ being meromorphic on the sphere $|s|=R$), and on the circle it is bounded above by
the constant $C_\xi R\log R$ (the growth bound \Cref{thm:xo_order}, with the $\log0=0$
convention covering any zeros on the circle). Hence by
\texttt{circleAverage\_mono\_on\_of\_le\_circle} and
$\operatorname{circleAverage}(\text{const})=\text{const}$
(\texttt{circleAverage\_const}),
\[
  N(R)\;\le\;C_\xi\,R\log R+\log2 .
\]
For $R\ge2$ one has $R\log R\ge2\log2>0$, so $\log2\le\tfrac12 R\log R$ and the right side
is $\le(C_\xi+\tfrac12)R\log R$. Take $C_N:=C_\xi+\tfrac12$.
\end{proof}

\section{The pointwise counting function}
\label{sec:zs_count}

\begin{definition}[Zero-counting function]
\label{def:zs_count}
\lean{riemannXiZeroCount}
\leanok
For $R\in\mathbb R$, the number of zeros of $\xi$ in the closed disc $|z|\le R$, counted
with multiplicity, is
\[
  n(R)\;:=\;\sum_{z}^{\flat} D_R(z)
        \;=\;\sum_{z}^{\flat}\ \texttt{MeromorphicOn.divisor riemannXi (closedBall 0\ R)}\ z .
\]
\end{definition}

\begin{lemma}[Pointwise counting bound]
\label{lem:zs_count_bound}
\lean{riemannXiZeroCount_le}
\leanok
\uses{lem:zs_jensen, def:zs_count}
There is a constant $C_n>0$ such that for all $R\ge2$,
\[
  n(R)\;\le\;C_n\,R\log R .
\]
\end{lemma}

\begin{proof}
\leanok
\uses{lem:zs_jensen, def:zs_count}
Compare $n(R)$ with $N(2R)$. Both finsums run over non-negative integer divisor values.
By \Cref{lem:zs_jensen} applied at radius $2R$ (note $2R\ge2$),
\[
  N(2R)\;=\;\sum_{z}^{\flat} D_{2R}(z)\,\log\!\frac{2R}{|z|}\;\le\;C_N\,(2R)\log(2R).
\]
Every summand on the left is non-negative: $D_{2R}\ge0$, and for $z$ in the support of
$D_{2R}$ one has $|z|\le2R$, hence $\log(2R/|z|)\ge0$. Dropping all summands with
$R<|z|\le2R$ (each $\ge0$) only decreases the sum, so
\[
  N(2R)\;\ge\;\sum_{|z|\le R}^{\flat} D_{2R}(z)\,\log\!\frac{2R}{|z|}
         \;=\;\sum_{|z|\le R}^{\flat} D_R(z)\,\log\!\frac{2R}{|z|},
\]
the last equality because $D_{2R}$ and $D_R$ agree on $\{|z|\le R\}$
(\texttt{toClosedBall\_eval\_within}; both equal $D(z)$ there). For $|z|\le R$ we have
$2R/|z|\ge2$, so $\log(2R/|z|)\ge\log2$; bounding each factor below by $\log2$ (all
weights $D_R(z)\ge0$) gives
\[
  N(2R)\;\ge\;(\log2)\sum_{|z|\le R}^{\flat} D_R(z)\;=\;(\log2)\,n(R).
\]
Combining, $n(R)\le N(2R)/\log2\le \dfrac{2C_N}{\log2}\,R\log(2R)$. Finally
$\log(2R)=\log2+\log R\le 2\log R$ for $R\ge2$ (since $\log2\le\log R$), so
$n(R)\le\dfrac{4C_N}{\log2}\,R\log R$. Take $C_n:=4C_N/\log2$.
\end{proof}

\section{Summability via dyadic decomposition}
\label{sec:zs_summable}

\begin{lemma}[Divisor-weighted summability]
\label{lem:zs_summable_divisor}
\lean{summable_riemannXiDivisor_norm_rpow}
\leanok
\uses{lem:zs_count_bound}
For every real $p>1$ the family
\[
  z\;\longmapsto\;\bigl(D(z)\bigr)\,\|z\|^{-p}\qquad(z\in\mathbb C)
\]
is summable (its support being the zeros of $\xi$, the summand vanishing off them and at
$z=0$).
\end{lemma}

\begin{proof}
\leanok
\uses{lem:zs_count_bound}
All summands are non-negative, so summability is equivalent to a uniform bound on the
finite partial sums; we bound the total by a convergent dyadic series. Group the zeros
by magnitude into the shells $A_k:=\{z:2^{k}\le|z|<2^{k+1}\}$ for $k\ge0$, together with
the inner disc $A_{-}:=\{z:|z|<1\}$. The inner disc contains finitely many zeros
($D$ has finite support on the compact disc $|z|\le1$,
\texttt{locallyFinsuppWithin.finiteSupport}), contributing a finite sum.

For the shells, on $A_k$ we have $\|z\|^{-p}\le 2^{-kp}$, and the total multiplicity in
$A_k$ is at most the count in the disc $|z|\le2^{k+1}$:
\[
  \sum_{z\in A_k}^{\flat} D(z)\;\le\;n\!\left(2^{k+1}\right)\;\le\;C_n\,2^{k+1}\,\log\!\left(2^{k+1}\right)
  \;=\;C_n\,(k+1)\,2^{k+1}\log2,
\]
using \Cref{lem:zs_count_bound} (and $2^{k+1}\ge2$). Hence the shell $A_k$ contributes at
most
\[
  2^{-kp}\sum_{z\in A_k}^{\flat} D(z)\;\le\;2C_n\log2\,\cdot\,(k+1)\,2^{-k(p-1)} .
\]
Since $p>1$ gives $0<2^{-(p-1)}<1$, the series $\sum_{k\ge0}(k+1)\,2^{-k(p-1)}$ converges
(\texttt{summable\_pow\_mul\_geometric\_of\_norm\_lt\_one} with $k$-power $1$ and ratio
$2^{-(p-1)}$). Adding the finite inner contribution, the total $\sum_z D(z)\|z\|^{-p}$ is
finite, which is the asserted summability.
\end{proof}

\section{The summability theorem and its downstream interface}
\label{sec:zs_main}

\begin{theorem}[Summability of the zeros of $\xi$]
\label{thm:zs_main}
\lean{summable_one_div_riemannZeta_zero_rpow}
\leanok
\uses{lem:zs_summable_divisor}
For every $\varepsilon>0$ the family
\[
  z\;\longmapsto\;\bigl(D(z)\bigr)\,\|z\|^{-(1+\varepsilon)}\qquad(z\in\mathbb C)
\]
is summable; that is, $\sum_\rho|\rho|^{-(1+\varepsilon)}<\infty$ over the zeros $\rho$ of
$\xi$ with multiplicity. In particular (taking $\varepsilon=1$),
$\sum_\rho|\rho|^{-2}<\infty$.
\end{theorem}

\begin{proof}
\leanok
\uses{lem:zs_summable_divisor}
Apply \Cref{lem:zs_summable_divisor} with $p:=1+\varepsilon>1$. The special case
$\varepsilon=1$ gives summability of $z\mapsto D(z)\|z\|^{-2}$.
\end{proof}

This is precisely \Cref{lem:lb_pf_summable} of \Cref{cha:lehman_bound}; the two
statements share the single Lean name
\texttt{summable\_one\_div\_riemannZeta\_zero\_rpow}. It remains to convert
the divisor-weighted form into the concrete enumeration consumed by the
canonical-product chapter.

\begin{lemma}[Downstream product hypothesis]
\label{lem:zs_product}
\lean{summable_inv_norm_sq_riemannXiZeros}
\leanok
\uses{thm:zs_main}
Let $\iota:=\bigl(z:\operatorname{supp}D\bigr)\times\operatorname{Fin}\,(D\,z).\mathrm{toNat}$
enumerate the zeros of $\xi$ with multiplicity, with $\rho:\iota\to\mathbb C$ the
first-projection map $\rho(z,j):=z$. Then $\sum_{i}\|\rho(i)^{-1}\|^{2}<\infty$; i.e.
\texttt{Summable (fun i ↦ ‖(ρ i)⁻¹‖²)}. Consequently the hypothesis of
\Cref{lem:hf_prod_summable}
(\lean{hasProdLocallyUniformlyOn_riemannXiCanonicalProduct}) is met for
this $\rho$, and the canonical product of \Cref{cha:hadamard} converges with $\xi$'s zeros.
\end{lemma}

\begin{proof}
\leanok
\uses{thm:zs_main}
For $i=(z,j)$ we have $\rho(i)=z\neq0$ (zeros of $\xi$ avoid $0$), so
$\|\rho(i)^{-1}\|^{2}=\|z\|^{-2}$, independent of $j$. Summing over the sigma index, the
inner sum over $\operatorname{Fin}\,(D\,z).\mathrm{toNat}$ multiplies by the multiplicity:
\[
  \sum_{i}\|\rho(i)^{-1}\|^{2}\;=\;\sum_{z\in\operatorname{supp}D}(D\,z).\mathrm{toNat}\;\|z\|^{-2}
  \;=\;\sum_{z\in\mathbb C}\bigl(D(z)\bigr)\|z\|^{-2},
\]
the last equality because $(D\,z).\mathrm{toNat}=D(z)$ on the support (where $D(z)\ge0$ is
a multiplicity) and both sides vanish off it. The right-hand family is summable by
\Cref{thm:zs_main} with $\varepsilon=1$; transport across the sigma reindexing via
\texttt{summable\_sigma\_of\_nonneg} (non-negative summands, finite inner fibres) and the
\texttt{finsum}/\texttt{tsum}-over-support identification yields the claim.
\end{proof}
